\section{Approximating maximum depth}
\label{sec:worst}

At a balanced state, one query immediately shrinks both child states.
At a flat state, there may be no such query, so the algorithm first
groups hypotheses by their answers to the hull-vertex tests and
then discovers which group contains the target. The two parts of
this work have different roles: the discovery schedule must finish
soon enough for each recursively constructed child code, while the
pattern class must be smaller than the parent. The analysis therefore
uses an optimal path only as a comparison schedule; the algorithm
uses the actual child-code heights when setting its deadlines.

Use the fixed rational count-balance threshold and coefficients
\[
 \delta=\frac{245122334}{10^9},\qquad
 \gamma=\frac{246496526}{10^8},\qquad
 \kappa=3+\gamma=5.46496526.
\]
They satisfy
\begin{equation}\label{eq:worst-parameter-conditions}
 \gamma\log_2\frac1{1-\delta}>1,\qquad
 \gamma\log_2\frac1\delta>5.
\end{equation}
These inequalities pay for balanced queries and flat-state decoding,
respectively. They can be checked with the rational logarithm bounds
in \eqref{eq:rational-log-bounds}.
The probabilities in the expected-depth model are not used by this
algorithm. Write $T_{\mathrm W}(S)$ for the code it constructs on
$S$; its actual height is $c_{\mathrm W}(T_{\mathrm W}(S))$.
In particular, $c_{\mathrm W}(T_{\mathrm W}(\{h\}))=0$.

\begin{symboltable}
$\delta,\ \gamma,\ \kappa$ & Count-balance threshold, coefficient of the logarithmic height credit, and approximation factor.\\
$s,\ D,\ H^*$ & Current state size, its optimum maximum depth, and target height $3D+\lfloor\gamma\log_2s\rfloor$.\\
$t_A,\ g_A$ & First supporting-block index on the comparison path, and actual first-positive position in the constructed schedule.\\
$e_A^{(4)},\ e_A^{(2)}$ & Endpoint-query allowances after a general or second-query positive exit.\\
$u_A,\ B_A,\ d_A$ & Child height plus endpoint allowance, remaining discovery budget under trial height $H$, and support-block deadline.\\
$\mathcal F,\ \mathcal S,\ \mathcal N$ & Classes positive on the first support, first positive on the second support, and negative on both.\\
\end{symboltable}

\subsection{Pattern classes and a height credit}

In a count-flat state $S$, let $s=|S|$ and choose the common witness
$p$ of Lemma~\ref{lem:helly}. Use the normalized hull, patterns,
and support blocks of Definitions~\ref{def:normalized-queries}
and~\ref{def:hull-pattern}.
Lemma~\ref{lem:hull-structure} makes $A_0$ a singleton when present.
Every nonzero class $A$ lies in a supplied minority hull cap, so
\[
 |A|<\delta s.
\]
Write $\Phi(s)=\lfloor\gamma\log_2s\rfloor$. The same integer
credit will pay for the maximum-depth constructions' endpoint exits.
\begin{lemma}[Count-potential credit]\label{lem:worst-floor-credit}
For a balanced child of size $a\le(1-\delta)s$,
$\Phi(s)-\Phi(a)\ge1$. For a flat-state nonzero class of size
$a<\delta s$, put $x=\log_2(s/a)$ and
$\ell=\log_2(1/\delta)$. Then
\begin{equation}\label{eq:worst-floor-credit}
 \Phi(s)-\Phi(a)\ge\lfloor\gamma x\rfloor
 \ge5+\lfloor x-\ell\rfloor.
\end{equation}
\end{lemma}
\begin{proof}
Use $\lfloor u\rfloor-\lfloor v\rfloor\ge\lfloor u-v\rfloor$.
The balanced claim follows from the first inequality in
\eqref{eq:worst-parameter-conditions}. In the flat case $x>\ell$,
and $\gamma\ge1$ and $\gamma\ell>5$ give
$\gamma x>5+(x-\ell)$. Taking floors proves the second claim.
\end{proof}

\begin{lemma}[Pathwise supporting-prefix credit]\label{lem:height-prefix}
Fix a proper maximum-depth-optimal tree for $S$, of height
$D=\OPT_{\mathrm W}(S)$, and follow $p$ along its path
$Q=(i_1,\ldots,i_q)$ to leaf $z$.
For a pattern class $A$, let $t_A$ be the first index whose
supporting hull vertices include a positive vertex of $A$;
put $t_A=q+1$ if there is no such index. Then
\[
 \OPT_{\mathrm W}(A)\le \OPT_{\mathrm W}(S)-(t_A-1).
\]
There is a comparison sequence of support blocks, each a singleton
or an adjacent pair of real hull vertices, covering every nonzero
pattern in its first $t_A$ blocks.
\end{lemma}
\begin{proof}
Before $t_A$, all supporting hull bits of each path query are zero
on the whole class. Convexity forces the original query's frozen
answer throughout $A$. Thus every member follows the same first
$t_A-1$ queries, and restriction to $A$ contracts all of them.
Every original path has length at most $D$, proving the height bound.
This is the height credit: the class can afford a longer search for
its pattern only to the extent that its common prefix has already
disappeared from the restricted optimal tree.

Replace each query of $Q$ by its block of real supporting
hull vertices. A class with $t_A\le q$ is covered in block $t_A$.
A class with no supporting hit follows the entire original path
and must be the singleton $\{z\}$. If that pattern is nonzero,
append one of its positive hull vertices as block $q+1$.
Repeated supports are harmless. If the pattern is zero, it is the
terminal core and needs no positive hit.
\end{proof}

\subsection{A pointwise pattern-decoding allowance}

Once a positive hull answer is known, the pattern can be recovered
by the endpoint method of Section~\ref{sec:endpoint-geometry}.
For the worst-case algorithm use the explicit midpoint-dyadic
alphabetic code in Lemma~\ref{lem:alphabetic-endpoints}, not a tree
chosen merely for minimum expected depth.
The code-design probabilities are uniform counts on $S$, conditioned
on the observed transcript; they do not represent randomization of
the algorithm.

\begin{lemma}[Per-class endpoint allowance]\label{lem:endpoint-height}
After the first positive answer, every hypothesis in a nonzero
class $A$, of size $a$, needs at most
\[
 e_A^{(4)}=\left\lfloor\log_2\frac{\delta s}{a}\right\rfloor+4
\]
additional queries to recover the full pattern.
If an adjacent hull vertex is already known to be negative, the
allowance improves to
\[
 e_A^{(2)}=\left\lfloor\log_2\frac{\delta s}{a}\right\rfloor+2.
\]
There is no additional query on the zero-pattern route once the
other patterns have been covered.
\end{lemma}
\begin{proof}
The opposite-edge cut in Section~\ref{sec:endpoint-geometry} makes
the positive pattern an ordinary interval without querying a
negative anchor. Its two endpoints are learned by supplied hull
thresholds.

Let $R$ be the compatible set after the first positive answer.
It lies in the positive hull cap, so $|R|<\delta s$.
The auxiliary uniform distribution on $S$ is uniform on $R$,
and all of $A$ remains compatible.
If $a_L$ members of $R$ have the target's left endpoint, the two
explicit alphabetic codes use at most
\[
 \log_2\frac{|R|}{a_L}+2
 +\log_2\frac{a_L}{|A|}+2
 =\log_2\frac{|R|}{a}+4
 <\log_2\frac{\delta s}{a}+4
\]
queries on that target. The second probability is conditioned on
the first endpoint. Their pair uniquely specifies its pattern.
Each of the two alphabetic searches contributes the additive $2$
in its pointwise code bound; multiplying the two conditional
probabilities makes their logarithmic terms telescope to
$\log_2(|R|/a)$.
The number of queries is an integer, so the displayed upper bound
implies the stated floor-based allowance.
If a negative vertex is adjacent to the positive one, cut at that
negative vertex and choose the orientation in which the positive
one is first. The positive interval must start there, so one
endpoint is fixed. A single alphabetic code learns the other
endpoint and uses at most
$\log_2(|R|/a)+2<\log_2(\delta s/a)+2$ queries.
This gives the second allowance. The actual negative answer is
the cut marker; no opposite-edge query or assumption about the
origin's location is required on this route.
In a hull of dimension at most one, the positive pattern is already
determined after discovery.
Singleton inference and cached answers only reduce the allowance.
\end{proof}

The next observation separates the local height test, which the
algorithm can check, from the optimal-path comparison used later.
\begin{lemma}[Local route-height test]\label{lem:worst-local-height}
Suppose each full pattern class $A$ has a stored identifying code
of actual height $h_A$. Discover a positive vertex using a fixed
ordering, truncated when every nonzero pattern is covered, and then
recover its full pattern. If $g_A$ is the first positive position,
the complete route through $A$ has height at most
\[
 g_A+h_A+e_A^{(4)},
\]
or $g_A+h_A+e_A^{(2)}$ when a negative neighbor is already cached.
If all these bounds are at most $H$, the entire tree has height at
most $H$, including the zero-pattern route when present.
\end{lemma}
\begin{proof}
The route consists of discovery, endpoint recovery, and the stored
child code; Lemma~\ref{lem:endpoint-height} bounds its middle part.
The all-negative route has length $\max_{A\ne A_0}g_A\le H$ and
ends at the zero singleton. An absent zero class creates no leaf.
For a nonsingleton state there is at least one nonzero pattern,
since $|A_0|\le1$.
\end{proof}

\subsection{Priority greedy with an enumerated adjacent pair}

The comparison path gives a deadline for every pattern class, but
the algorithm does not know that path. It must produce one ordering
that meets all deadlines at once. Enumerating the first adjacent
pair fixes the circular cut; priority weights then make ordinary
maximum-remaining-weight greedy behave like greedy for every
deadline prefix, without requiring the comparison blocks as input.

The indices $t_A,d_A$ count support blocks; $g_A$ counts actual
vertex queries. The scheduler translates a feasible block deadline
into a bound on that actual first-positive position.

\begin{lemma}[Adjacent-block schedule]\label{lem:deadline-cover}
Fix two distinct adjacent hull vertices $v,w$ and consider the
pattern classes negative on both. Assign each such class an integer
block budget $d_A\ge2$. Suppose a comparison sequence of support
blocks starts with $\{v,w\}$, every other block is empty, a singleton,
or an adjacent pair, and it hits $A$ in its first $d_A$ blocks.

Force $v,w$ along the negative discovery route. Give a class of
budget $d^{(j)}$ the auxiliary integer weight $M^{r-j}$, where
$d^{(1)}<\cdots<d^{(r)}$ are the distinct budgets and $M=s+1$.
Then complete discovery by maximum-remaining-weight greedy.
In this one ordering every surviving class satisfies
\[
 g_A\le3d_A-2.
\]
The schedule and its first-hit times are computable in polynomial
time without knowing the comparison blocks.
\end{lemma}
\begin{proof}
Fix a budget $d=d^{(j)}$ and restrict demands to the prefix
$\mathcal P_j=\{A:d_A\le d\}$. Every score from these demands
is an integer multiple of $M^{r-j}$. The total weight of all later
demands is less than $M^{r-j}$: there are at most $s$ classes and
each later weight is at most $M^{r-j-1}$, with $s<M$.
When there are no later demands this weight is zero.

The weight of one prefix hit therefore outweighs every possible
contribution from later-deadline classes. Thus later demands cannot
reverse a strict comparison of prefix scores. After the forced pair,
until the prefix is covered,
each global greedy choice
is a weighted greedy choice for its remaining prefix demands.
It hits at least one of them; a zero prefix score cannot beat a
positive prefix score. Ties may use later weights but are legitimate
greedy ties for the prefix.

Cut at $v$ and delete $w$. The remaining demands are ordinary
intervals. Because the removed vertices were adjacent, every
comparison pair avoiding them is still consecutive in this linear
candidate order; a block meeting them becomes at most a singleton.
Let $P$ be the union of the first $d$ comparison blocks with $v,w$
removed, and let $c(P)$ count its runs of consecutive candidates.
The first block contributes nothing. Each later block contributes
at most three to size plus run count: an adjacent pair has size two
and one run, while a singleton has size one and one run. Both
quantities are subadditive under union, so
\[
 |P|+c(P)\le3(d-1).
\]
This set covers the surviving prefix. For any nonempty surviving
prefix, its classes are negative on $v,w$ and therefore must be hit
in a later comparison block, so $P$ is nonempty.
If it is nonempty,
Lemma~\ref{lem:greedy-cover-length} bounds the number of additional
greedy choices by $|P|+c(P)-1\le3d-4$.
Including the forced pair gives $3d-2$.
The priority argument applies to every prefix in the same ordering.

The auxiliary weights have $O(s\log s)$ bits, and greedy only sums
and compares them using the explicit class-to-vertex incidence
table. The block budgets are internal scheduling quantities, not
additional constraints in the identification problem.
\end{proof}

\begin{algorithm}[H]
\caption{\PriorityTrial: one height trial with stored child codes}
\label{alg:priority-trial}
\KwIn{Flat state's hull and classes; child heights $h_A$; allowances $e_A^{(4)},e_A^{(2)}$; target $H$.}
\KwOut{A passing discovery order, or infeasible.}
\ForEach{oriented adjacent real pair $(v,w)$}{
  Form $\mathcal F,\mathcal S,\mathcal N$ by their first hit on $(v,w)$\;
  \If{some $A\in\mathcal F$ has $1+h_A+e_A^{(4)}>H$, or
  some $A\in\mathcal S$ has $2+h_A+e_A^{(2)}>H$}{\Continue\;}
  \ForEach{$A\in\mathcal N$}{
    $u_A\gets h_A+e_A^{(4)}$;
    $B_A\gets H-u_A$;
    $d_A\gets\lfloor(B_A+2)/3\rfloor$\;
  }
  \If{some $B_A<4$}{\Continue\;}
  Sort distinct budgets as $d^{(1)}<\cdots<d^{(r)}$\;
  Give each deadline-$d^{(j)}$ class weight $(|S|+1)^{r-j}$\;
  Force $v,w$ and complete by maximum-remaining-weight greedy;
  truncate when all nonzero classes are hit\;
  \If{every $A\in\mathcal N$ has actual $g_A+u_A\le H$}{\KwRet this order\;}
}
\KwRet infeasible\;
\end{algorithm}


\subsection{The recursive maximum-depth algorithm}

Algorithms~\ref{alg:priority-trial} and~\ref{alg:worst-tree} give the
explicit construction. At a balanced state query a balanced line and
recurse. At a flat state build every pattern child \emph{before}
scheduling; the trial uses its actual height, not its optimum. For an
oriented initial pair $(v,w)$, the first-positive classes satisfy the
direct one-query and two-query height checks. For classes negative on
both, put
\[
 u_A=c_{\mathrm W}(T_{\mathrm W}(A))+e_A^{(4)},\qquad
 B_A=H-u_A,\qquad d_A=\lfloor(B_A+2)/3\rfloor.
\]
Reject if $B_A<4$, then use the priority weights of
Lemma~\ref{lem:deadline-cover} and accept only if each actual first hit
satisfies $g_A+u_A\le H$. Trial heights range from one through
$\lfloor\kappa(s-1)\rfloor$.

After a positive answer, Algorithm~\ref{alg:endpoint-code} learns the
full pattern and the stored child identifies its hypothesis. The
second-query exit has a cached negative neighbor and uses only one
endpoint code. Stop discovery as soon as all nonzero patterns are
covered; its all-negative route identifies $A_0$ when present and is
omitted otherwise. Low-dimensional hulls use their extreme tests
directly. Every nonsingleton pattern child is smaller, so the recursion
is well founded.

Lemma~\ref{lem:worst-local-height} makes a passing trial a complete
height certificate. Its first two exits have discovery costs one
and two, respectively; the latter has a cached negative neighbor.
Every remaining class uses the general allowance. The checks use
the explicit schedule's actual first-hit times.

\begin{algorithm}[H]
\caption{\WorstTree: the priority maximum-depth baseline}
\label{alg:worst-tree}
\KwIn{State $S$ and the rational parameters of Section~\ref{sec:worst}.}
\KwOut{A supplied-query identification tree.}
\If{$|S|\le1$}{\KwRet its leaf\;}
\If{a count-balanced supplied line $\ell$ exists}{
  \KwRet $\ell$ with recursively constructed children $\WorstTree(S_\ell^-)$ and $\WorstTree(S_\ell^+)$\;
}
Compute $\HullData(S)$ and recursively build one child $T_A$ per pattern class\;
\If{the hull has at most two real extreme vertices}{
  \KwRet their direct pattern stage with the stored children attached\;
}
Compute $e_A^{(4)},e_A^{(2)}$ from Lemma~\ref{lem:endpoint-height}\;
\For{$H=1,\ldots,\lfloor(3+\gamma)(|S|-1)\rfloor$}{
  $Q\gets\PriorityTrial(S,(c_W(T_A)),(e_A^{(4)},e_A^{(2)}),H)$\;
  \If{$Q$ is feasible}{
    Build its first-positive exits, use pointwise $\EndpointCode$ at each exit,
    and attach the stored class children\;
    End the all-negative route at the zero singleton when present\;
    \KwRet this tree\;
  }
}
\end{algorithm}


\begin{theorem}[Constant deterministic maximum-depth approximation]
\label{thm:worst}
The preceding polynomial-time deterministic algorithm satisfies
\[
 c_{\mathrm W}(T_{\mathrm W}(S))\le
 3\OPT_{\mathrm W}(S)+\left\lfloor\gamma\log_2|S|\right\rfloor
 \le\kappa\OPT_{\mathrm W}(S),\qquad \kappa=5.46496526<5.464966.
\]
It uses only supplied line queries.
\end{theorem}
\begin{proof}
Induct on $s=|S|$, with a zero-cost singleton, and write
$D=\OPT_{\mathrm W}(S)$ and $H^*=3D+\lfloor\gamma\log_2s\rfloor$.
A balanced child has size $a\le(1-\delta)s$ and optimum at most $D$,
so Lemma~\ref{lem:worst-floor-credit} pays for its root query.

In a flat state put $L_A=\log_2(s/|A|)$ and
$\ell=\log_2(1/\delta)$.

Choose an enumerated initial pair containing the first optimal-witness
query's supports; pad a singleton support with an adjacent real
neighbor. This only adds coverage. The inductive bound and
\eqref{eq:worst-floor-credit} leave at least one query after the
general endpoint allowance, and at least three after the singly pinned
allowance. Thus the first and second positive exits pass their direct
height checks, including classes hit by the padding.

For a class $A$ negative on both first supports, its comparison index
has $t_A\ge2$. Lemma~\ref{lem:height-prefix} and induction give
\[
\begin{aligned}
 B_A=H^*-u_A
 &\ge3(t_A-1)+\lfloor\gamma L_A\rfloor
                    -\lfloor L_A-\ell\rfloor-4\\
 &\ge3t_A-2\ge4.
\end{aligned}
\]
Hence $d_A=\lfloor(B_A+2)/3\rfloor\ge t_A$. The padded first block,
the original supporting blocks, and the terminal nonzero singleton
block when needed satisfy Lemma~\ref{lem:deadline-cover}. Its single
priority ordering therefore obeys
\[
 g_A\le3d_A-2
 =3\lfloor(B_A+2)/3\rfloor-2\le B_A.
\]
The last inequality follows by writing $B_A+2=3d_A+r$ with
$r\in\{0,1,2\}$. All actual-height checks pass. Truncating at the
last positive first hit also puts the zero singleton within $H^*$.
A one-dimensional hull needs at most two extreme tests; its count
potential drop pays for these directly.

Thus the least passing trial has height at most $H^*$. Since
$\log_2s\le D\le s-1$, this integer trial lies in the prescribed
range and $H^*\le(3+\gamma)D$. No optimum is queried by the algorithm.
\end{proof}

The bound controls the longest path of one deterministic output
tree, not an expected query count over random trees.
There are at most $n-1$ nonterminal stages. A flat stage tries
$O(s)$ integer bounds, each trying at most $2|V|$ oriented adjacent
pairs and a weighted greedy completion.
The endpoint code-design probabilities are counts, so
the explicit alphabetic codewords have $O(\log s)$ bits.
Geometry, child-height evaluation, schedule checks, and decoding-code
construction all have polynomial bit complexity.
The proper complete output tree has $n-1$ internal nodes.

The two parameter constraints explain the optimized threshold.
For this induction, the least possible coefficient is
\[
 \gamma(\delta)=
 \max\left\{\frac1{\log_2(1/(1-\delta))},
             \frac5{\log_2(1/\delta)}\right\}.
\]
The first term decreases with $\delta$ and the second increases,
so their maximum is minimized where they are equal. Solving this
equation gives $\delta=(1-\delta)^5$, at approximately
$0.2451223337533$.
The chosen rational parameters give a factor within $5\cdot10^{-9}$
of the resulting minimum $5.4649652555$.
The first pair also serves two purposes: it removes one support
block from the greedy coverage bound, and its second-query exit is
singly pinned, so only one endpoint code is needed.
